\chapter{Stochastic Differential Equations}

\begin{center}\itshape
A particle pushed by a force and kicked by noise, a price that grows and trembles, eigenvalues that wander and repel:\\
each is described by an equation in which $\mathrm dB$ appears next to $\mathrm dt$.\\
This chapter solves such equations, and finds that some solutions refuse to do what they seem able to do.
\end{center}

In 1908, three years after Einstein's paper on Brownian motion, Paul Langevin presented to the Paris Academy of Sciences a derivation of Einstein's formula for the mean square displacement which was,
in his words, ``infinitely more simple'' (in the English translation of Lemons and Gythiel, 1997). Einstein had written a diffusion equation for the density of particles. Langevin wrote Newton's law for
a single particle: mass times acceleration equals a friction force proportional to the velocity, plus a ``complementary force'' that is irregular, as likely to push one way as the other, and
independent of the position. He then averaged, and the formula came out in a few lines. What his complementary force was, as a mathematical object, nobody could say. The velocity of his particle
satisfied, in modern notation, $\mathrm dV=-\theta V\dd t+\sigma\dd B$.

In 1930, George Uhlenbeck and Leonard Ornstein computed the law of Langevin's velocity: it is Gaussian, it forgets its initial value exponentially fast, and it settles into an equilibrium. Their paper
in the \emph{Physical Review} gave the process its name. Twenty years later, It\^o's integral gave Langevin's equation a meaning, and It\^o proved in 1951 that such equations have exactly one solution
when their coefficients are Lipschitz.

The next equation came from finance. In 1973 Fischer Black and Myron Scholes, and independently Robert Merton, priced an option on a stock whose price follows $\mathrm dS=\mu S\dd t+\sigma S\dd B$. Their
argument is an application of It\^o's formula: a portfolio of the stock and a bond, readjusted continuously, reproduces the option exactly, so the option has a single fair price, in which the
expected return $\mu$ of the stock does not even appear. The paper of Black and Scholes was first rejected by the \emph{Journal of Political Economy}, which published it after a second look, as Black
later recalled (1989). Merton and Scholes received the Nobel prize in economics in 1997; Black had died in 1995.

The last equation of this chapter is Freeman Dyson's. In 1962 he let the entries of a Hermitian matrix perform independent Brownian motions and asked how its eigenvalues move. They perform Brownian
motions too, but they repel each other with a force inversely proportional to their distance, a two-dimensional Coulomb gas on a line. Dyson's paper was called \emph{A Brownian-motion model for the
eigenvalues of a random matrix}; Henry McKean found, a few years later, a short argument that the eigenvalues never collide. That argument, a Lyapunov function and It\^o's formula, ends the chapter
before a theorem on the numerical solution of all these equations.

\begin{reminder}[from Chapters 3 to 5]
$B$ is an $m$-dimensional $(\mathcal F_t)$-Brownian motion, the filtration satisfying the usual conditions. For progressive $H$ with $\mathbb E\int_0^tH_s^2\dd s<\infty$, $\int_0^tH\dd B^j$ is a
square-integrable martingale with $\mathbb E(\int_0^tH\dd B^j)^2=\mathbb E\int_0^tH^2\dd s$, and $\langle H\cdot B^j,K\cdot B^k\rangle=\delta_{jk}\int HK\dd s$ (Theorem~\ref{thm:itoloc}, Proposition~\ref{prop:itoprops}). Doob's
inequality $\mathbb E\sup_{s\le t}M_s^2\le4\mathbb EM_t^2$ (Theorem~\ref{thm:doobcont}). It\^o's formula (Theorem~\ref{thm:itoformula}), also for functions that are $C^1$ in time and $C^2$ in space (the paragraph
before Theorem~\ref{thm:feynmankac}). L\'evy's characterisation (Theorem~\ref{thm:levychar}).
\end{reminder}

\section{Existence and uniqueness}

Let $b:\R^d\to\R^d$ and $\sigma:\R^d\to\R^{d\times m}$ be measurable. For a $d\times m$ matrix $A$ we write $|A|$ for its Euclidean (Frobenius) norm, and for a progressive $\R^{d\times m}$-valued $H$,
$\int_0^tH\dd B$ is the vector whose $i$-th coordinate is $\sum_j\int_0^tH_{ij}\dd B^j$. The reminder gives the isometry
\begin{equation}\label{eq:isomvec}
\mathbb E\Bigl|\int_0^tH\dd B\Bigr|^2=\mathbb E\int_0^t|H_s|^2\dd s ,
\end{equation}
since the cross terms vanish: for $j\ne j'$,
\[
\mathbb E\Bigl[\int H_{ij}\dd B^j\int H_{ij'}\dd B^{j'}\Bigr]=\mathbb E\int H_{ij}H_{ij'}\dd\langle B^j,B^{j'}\rangle=0.
\]

\begin{definition}\label{def:strongsol}
Let $\xi$ be an $\mathcal F_0$-measurable random vector. A \emph{(strong) solution} of
\[
\dd X_t=b(X_t)\dd t+\sigma(X_t)\dd B_t,\qquad X_0=\xi,
\]
is a continuous adapted process $X$ in $\R^d$ such that, almost surely, $\int_0^t(|b(X_s)|+|\sigma(X_s)|^2)\dd s<\infty$ and $X_t=\xi+\int_0^tb(X_s)\dd s+\int_0^t\sigma(X_s)\dd B_s$ for all $t\ge0$.
\end{definition}

\begin{lemma}[Gronwall]\label{lem:gronwall}
Let $\varphi$ be bounded, measurable and non-negative on $[0,T]$, with $\varphi(t)\le a+L\int_0^t\varphi(s)\dd s$ for all $t\le T$, where $a,L\ge0$. Then $\varphi(t)\le ae^{Lt}$ for $t\le T$.
\end{lemma}

\begin{proof}
Substituting the inequality into itself $n$ times gives
\[
\varphi(t)\le a\sum_{k<n}\frac{(Lt)^k}{k!}+L^n\int_{0\le s_n\le\dots\le s_1\le t}\varphi(s_n)\dd s_n\cdots\dd s_1 ,
\]
and the last term is at most $\sup\varphi\cdot\frac{(Lt)^n}{n!}\to0$.
\end{proof}

\begin{theorem}[It\^o, 1951]\label{thm:sdeexist}
Assume that $b$ and $\sigma$ are Lipschitz: $|b(x)-b(y)|+|\sigma(x)-\sigma(y)|\le K|x-y|$. For every $\mathcal F_0$-measurable $\xi$ with $\mathbb E|\xi|^2<\infty$ there is a solution $X$, unique up to
indistinguishability, and for every $T$ there is $C_T$, depending only on $T$, $K$, $|b(0)|$, $|\sigma(0)|$, with
\[
\mathbb E\sup_{t\le T}|X_t|^2\le C_T(1+\mathbb E|\xi|^2).
\]
\end{theorem}

\begin{proof}
Fix $T$ and let $\mathcal S_T$ be the space of continuous adapted processes on $[0,T]$ with $\|X\|^2=\mathbb E\sup_{t\le T}|X_t|^2<\infty$; it is complete (a Cauchy sequence has a subsequence converging
uniformly almost surely, as in Theorem~\ref{thm:L2cont}). For $X\in\mathcal S_T$ put
\[
\Phi(X)_t=\xi+\int_0^tb(X_s)\dd s+\int_0^t\sigma(X_s)\dd B_s.
\]
Since $|b(x)|\le|b(0)|+K|x|$ and similarly for $\sigma$, the integrands are square integrable, and by Cauchy--Schwarz for the first integral, Doob's inequality and \eqref{eq:isomvec} for the second,
$\Phi(X)\in\mathcal S_T$. Write $D_t(X,Y)=\mathbb E\sup_{s\le t}|X_s-Y_s|^2$. The same two inequalities give, for $X,Y\in\mathcal S_T$ and $t\le T$, with $L=2(T+4)K^2$,
\begin{equation}\label{eq:picard}
D_t\bigl(\Phi(X),\Phi(Y)\bigr)\le L\int_0^tD_s(X,Y)\dd s.
\end{equation} Start from $X^0=\xi$ and put $X^{n+1}=\Phi(X^n)$. By induction from \eqref{eq:picard}, $D_t(X^{n+1},X^n)\le\frac{(Lt)^n}{n!}\|X^1-X^0\|^2$, whose square roots are
summable. So $(X^n)$ is Cauchy in $\mathcal S_T$ and converges to some $X$; by \eqref{eq:picard} $\Phi$ is continuous on $\mathcal S_T$, so $X=\Phi(X)$: $X$ solves the equation on $[0,T]$.

\emph{Uniqueness.} Let $X$ and $Y$ be two solutions on $[0,T]$ (not assumed square integrable) and $\tau_k=\inf\{t:|X_t|\vee|Y_t|\ge k\}$. On $\{s\le\tau_k\}$ we have $|X_s-Y_s|\le|X_{s\wedge\tau_k}-Y_{s\wedge\tau_k}|$, and
the stopped integrals are integrals of $\mathbf 1_{\{s\le\tau_k\}}$ times the integrands (Proposition~\ref{prop:itoprops}(c)), so the proof of \eqref{eq:picard} applies to
$\varphi(t)=\mathbb E\sup_{s\le t}|X_{s\wedge\tau_k}-Y_{s\wedge\tau_k}|^2\mathbf 1_{\{\tau_k>0\}}$, which is bounded by $4k^2$:
$\varphi(t)\le L\int_0^t\varphi$. Gronwall gives $\varphi=0$, so $X=Y$ on $[0,\tau_k]$, and $\tau_k\to\infty$.

\emph{The bound.} $\mathbb E\sup_{s\le t}|X_s|^2\le3\mathbb E|\xi|^2+3T\,\mathbb E\int_0^t|b(X_s)|^2\dd s+12\,\mathbb E\int_0^t|\sigma(X_s)|^2\dd s$, and $|b(x)|^2\le2|b(0)|^2+2K^2|x|^2$; Gronwall applied to
$\varphi(t)=\mathbb E\sup_{s\le t}|X_s|^2$, which is finite, gives the bound.

Finally the solutions on $[0,N]$, $N=1,2,\dots$, agree on their common intervals by uniqueness, and define a solution on $[0,\infty[$.
\end{proof}

The proof is Picard's iteration for ordinary differential equations, with Doob's inequality in place of the supremum bound of the integral. The same Gronwall argument compares two solutions with
different initial values.

\begin{proposition}\label{prop:sdeflow}
Under the hypotheses of Theorem~\ref{thm:sdeexist}, if $X$ and $Y$ are the solutions from $\xi$ and $\eta$, then $\mathbb E\sup_{t\le T}|X_t-Y_t|^2\le3e^{LT}\,\mathbb E|\xi-\eta|^2$, with $L=6(T+4)K^2$.
\end{proposition}

\begin{proof}
$X-Y=\xi-\eta+\int(b(X)-b(Y))\dd s+\int(\sigma(X)-\sigma(Y))\dd B$, so as in \eqref{eq:picard}, with $(a+b+c)^2\le3(a^2+b^2+c^2)$, $\varphi(t)=\mathbb E\sup_{s\le t}|X_s-Y_s|^2$ satisfies
$\varphi(t)\le3\mathbb E|\xi-\eta|^2+L\int_0^t\varphi$. Apply Gronwall.
\end{proof}

For $x\in\R^d$ let $X^x$ be the solution from $x$, and define $P_tf(x)=\mathbb Ef(X^x_t)$ for bounded measurable $f$. The solution forgets the past: at time $t$ it starts afresh from where it is.

\begin{theorem}[Markov property]\label{thm:sdemarkov}
Under the hypotheses of Theorem~\ref{thm:sdeexist}, $x\mapsto P_sf(x)$ is measurable, and for all $x$, $s,t\ge0$ and bounded measurable $f$,
\[
\mathbb E\bigl[f(X^x_{t+s})\mid \mathcal F_t\bigr]=P_sf(X^x_t)\qquad\text{almost surely.}
\]
\end{theorem}

\begin{proof}
Let $B'_u=B_{t+u}-B_t$, a Brownian motion for the filtration $\mathcal G_u=\mathcal F_{t+u}$, independent of $\mathcal G_0=\mathcal F_t$, and $Y_u=X^x_{t+u}$. Then
$Y_u=X^x_t+\int_0^ub(Y_v)\dd v+\int_0^u\sigma(Y_v)\dd B'_v$: for elementary integrands the shift of the stochastic integral is clear, and both sides are continuous for the isometry. So $Y$ is the solution
driven by $B'$, for the filtration $(\mathcal G_u)$, from the $\mathcal G_0$-measurable initial value $\xi=X^x_t$.

For $y\in\R^d$ let $Z^y$ be the solution from $y$ driven by $B'$ in the (augmented) filtration of $B'$. The Picard iterates are built from $B'$ alone, so $Z^y$ is measurable with respect to $\sigma(B')$, which is
independent of $\mathcal F_t$, and $Z^y$ has the law of $X^y$, since the law of each iterate depends only on the law of the driving Brownian motion. A stochastic integral with respect to $B'$ of a process
adapted to the smaller filtration is the same in both filtrations (again clear for elementary integrands), so $Z^y$ also solves the equation for $(\mathcal G_u)$.

Let $\xi_k$ be $\mathcal F_t$-measurable with finitely many values $y_1,\dots,y_r$ and $|\xi_k-\xi|\le1/k$. Then $\hat Y^k=\sum_i\mathbf 1_{\{\xi_k=y_i\}}Z^{y_i}$ solves the equation from $\xi_k$ for $(\mathcal G_u)$
(the indicators are $\mathcal G_0$-measurable and pass through the integrals), so by Proposition~\ref{prop:sdeflow} $\mathbb E\sup_{u\le s}|\hat Y^k_u-Y_u|^2\le Ck^{-2}$. For $f$ bounded and Lipschitz,
independence gives
\[
\mathbb E[f(\hat Y^k_s)\mid\mathcal F_t]=\sum_i\mathbf 1_{\{\xi_k=y_i\}}\mathbb E f(Z^{y_i}_s)=P_sf(\xi_k).
\]
As $k\to\infty$, the left side tends to $\mathbb E[f(Y_s)\mid\mathcal F_t]$ in $L^1$, and the right side to $P_sf(\xi)$, because $|P_sf(y)-P_sf(z)|\le\mathrm{Lip}(f)\,\mathbb E|X^y_s-X^z_s|\le C|y-z|$ by
Proposition~\ref{prop:sdeflow}. This proves the identity for bounded Lipschitz $f$, for which $P_sf$ is continuous. For a closed set $F$, $f_n=\max(0,1-n\,d(\cdot,F))$ decreases to $\mathbf 1_F$, and monotone convergence
extends both the measurability and the identity to $\mathbf 1_F$; the class of Borel sets $A$ for which they hold for $\mathbf 1_A$ is a $\lambda$-system containing the $\pi$-system of closed sets, hence all Borel sets,
and linearity and monotone convergence give all bounded measurable $f$.
\end{proof}

Lipschitz conditions are not a technicality. The equation $\dd X=3X^{1/3}\dd t+3X^{2/3}\dd B$, $X_0=0$, has the two solutions $X=0$ and $X=B^3$ (Exercise~\ref{ex:nonunique}). In dimension one, however,
the diffusion coefficient may be much rougher than Lipschitz.

\begin{theorem}[Yamada and Watanabe, 1971]\label{thm:yamada}
Let $d=m=1$, $b$ Lipschitz and $\sigma$ such that $|\sigma(x)-\sigma(y)|\le K\sqrt{|x-y|}$. If $X$ and $Y$ are two solutions with $X_0=Y_0$, driven by the same Brownian motion, then $X$ and $Y$ are indistinguishable.
\end{theorem}

\begin{proof}
The idea is to apply It\^o's formula to $|X-Y|$, smoothed so that the second derivative is of order $1/(n|x|)$ near $0$, which exactly compensates the $|x|$ coming from $|\sigma(X)-\sigma(Y)|^2\le K^2|X-Y|$.

Let $a_0=1$ and $a_n=a_{n-1}e^{-n}$, so that $\int_{a_n}^{a_{n-1}}\frac{\mathrm du}u=n$. Choose continuous $\rho_n\ge0$ supported in $]a_n,a_{n-1}[$ with $\rho_n(u)\le\frac2{nu}$ and $\int\rho_n=1$, and put
$\psi_n(x)=\int_0^{|x|}\int_0^y\rho_n(u)\dd u\dd y$. Then $\psi_n\in C^2(\R)$ (it vanishes near $0$), $|\psi_n'|\le1$, $\psi_n''(x)=\rho_n(|x|)\le\frac2{n|x|}$, and $\psi_n(x)\uparrow|x|$ as $n\to\infty$,
since $\psi'_n(y)=1$ for $y\ge a_{n-1}$.

Let $\Delta=X-Y$ and $\tau_k=\inf\{t:|X_t|\vee|Y_t|\ge k\}$. By It\^o's formula,
with $\delta b_s=b(X_s)-b(Y_s)$ and $\delta\sigma_s=\sigma(X_s)-\sigma(Y_s)$,
\[
\psi_n(\Delta_t)=\int_0^t\psi_n'(\Delta_s)\,\delta b_s\dd s+\int_0^t\psi_n'(\Delta_s)\,\delta\sigma_s\dd B_s+\frac12\int_0^t\psi_n''(\Delta_s)\,\delta\sigma_s^2\dd s.
\]
Stopped at $\tau_k$, the stochastic integral has a bounded integrand and mean $0$. The last term is at most $\frac12\int_0^t\frac2{n|\Delta_s|}K^2|\Delta_s|\dd s=\frac{K^2t}n$. Hence
\[
\mathbb E\psi_n(\Delta_{t\wedge\tau_k})\le K\,\mathbb E\int_0^{t\wedge\tau_k}|\Delta_s|\dd s+\frac{K^2t}n\le K\int_0^t\mathbb E|\Delta_{s\wedge\tau_k}|\dd s+\frac{K^2t}n.
\]
Let $n\to\infty$ (monotone convergence): $\varphi(t)=\mathbb E|\Delta_{t\wedge\tau_k}|\le2k$ satisfies $\varphi(t)\le K\int_0^t\varphi$, so $\varphi=0$ by Gronwall. Thus $X_t=Y_t$ almost surely on $\{t\le\tau_k\}$ for each $t$ and
$k$; let $k\to\infty$ and use continuity.
\end{proof}

\begin{figure}[ht]
\centering
\fig{ch06-ou}
\caption{Three paths of the Ornstein--Uhlenbeck process $\mathrm dX=-X\dd t+\dd B$, started at $3$, $0$ and $-3$. They forget their starting points at the exponential rate $e^{-t}$ and then fluctuate in the
stationary law $N(0,\frac12)$; the dashed lines are at $\pm2/\sqrt2$, two standard deviations.}
\label{fig:ou}
\end{figure}

\section{Linear equations: Langevin and Black--Scholes}

\begin{proposition}[Ornstein--Uhlenbeck, 1930]\label{prop:ou}
Let $\theta,\sigma>0$. The solution of $\dd X=-\theta X\dd t+\sigma\dd B$, $X_0=x$, is
\[
X_t=e^{-\theta t}x+\sigma\int_0^te^{-\theta(t-s)}\dd B_s ,
\]
a Gaussian process with $X_t\sim N\bigl(e^{-\theta t}x,\frac{\sigma^2}{2\theta}(1-e^{-2\theta t})\bigr)$. As $t\to\infty$, $X_t$ converges in law to $N(0,\frac{\sigma^2}{2\theta})$, and if $X_0$ is independent of $B$ with this law, then
$X_t$ has this law for every $t$.
\end{proposition}

\begin{proof}
It\^o's formula for $f(t,x)=e^{\theta t}x$ gives $\dd(e^{\theta t}X_t)=\theta e^{\theta t}X_t\dd t+e^{\theta t}\dd X_t=\sigma e^{\theta t}\dd B_t$, whence the formula; conversely the formula defines a solution by the same
computation, and the solution is unique (Theorem~\ref{thm:sdeexist}). The integral of the deterministic function $e^{-\theta(t-s)}$ is a Wiener integral (Exercise~\ref{ex:wienerito}), Gaussian, centred, with
variance $\sigma^2\int_0^te^{-2\theta(t-s)}\dd s=\frac{\sigma^2}{2\theta}(1-e^{-2\theta t})$, and jointly Gaussian at several times. The limit in law follows from the characteristic functions. If $X_0\sim N(0,\frac{\sigma^2}{2\theta})$ is
independent of $B$, then $X_t$ is the sum of the independent Gaussians $e^{-\theta t}X_0$ and $\sigma\int_0^te^{-\theta(t-s)}\dd B_s$, with variance $\frac{\sigma^2}{2\theta}e^{-2\theta t}+\frac{\sigma^2}{2\theta}(1-e^{-2\theta t})$.
\end{proof}

With $\sigma=\sqrt2$, $\theta=1$ and $X_0\sim N(0,1)$ this is the stationary process of covariance $e^{-|t-s|}$ met in Chapter~1 (Exercise~\ref{ex:ou}). Langevin's velocity has a position $\int_0^tX_s\dd s$; over long times it behaves like
a Brownian motion with variance $\frac{\sigma^2}{\theta^2}t$ (Exercise~\ref{ex:langevin}), which is Einstein's diffusion recovered from Newton's law.

\begin{proposition}[geometric Brownian motion]\label{prop:gbm}
Let $\mu\in\R$, $\sigma>0$, $s>0$. The solution of $\dd S=\mu S\dd t+\sigma S\dd B$, $S_0=s$, is $S_t=s\exp\bigl(\sigma B_t+(\mu-\frac{\sigma^2}2)t\bigr)$.
\end{proposition}

\begin{proof}
It\^o's formula for $f(t,y)=s\exp(\sigma y+(\mu-\frac{\sigma^2}2)t)$ gives $\dd f(t,B_t)=f\cdot\bigl((\mu-\frac{\sigma^2}2)\dd t+\sigma\dd B_t+\frac{\sigma^2}2\dd t\bigr)=\mu S\dd t+\sigma S\dd B$; the coefficients are linear, so
the solution is unique.
\end{proof}

The correction $-\frac{\sigma^2}2$ is It\^o's: the median of $S_t$ grows at rate $\mu-\frac{\sigma^2}2$, the mean at rate $\mu$. If $0<\mu<\frac{\sigma^2}2$, the mean tends to infinity while almost every path tends
to $0$ (Exercise~\ref{ex:gbm}).

\medskip
\emph{The Black--Scholes model.} A market has a bond worth $e^{rt}$ at time $t$ (interest rate $r\ge0$) and a stock with price $S$ as in Proposition~\ref{prop:gbm}. A \emph{European call} with strike $K>0$ and
maturity $T$ pays $(S_T-K)^+$ at time $T$: the right to buy the stock at price $K$. A \emph{strategy} holds $a_t$ shares and $c_t$ bonds, with $a,c$ progressive and locally bounded on $[0,T[$; its value is
$V_t=a_tS_t+c_te^{rt}$, and it is \emph{self-financing} if
\[
V_t=V_0+\int_0^ta_u\dd S_u+\int_0^tc_u\,re^{ru}\dd u\qquad(t<T),
\]
that is, if the value changes only through the price movements, without money added or withdrawn. Let $\Phi$ and $\varphi$ be the standard normal distribution function and density, $\tau=T-t$, and
\[
d_\pm(t,s)=\frac{\log(s/K)+(r\pm\frac{\sigma^2}2)\tau}{\sigma\sqrt\tau},\qquad C(t,s)=s\,\Phi(d_+)-Ke^{-r\tau}\Phi(d_-)\quad(t<T),
\]
and $C(T,s)=(s-K)^+$.

\begin{theorem}[Black, Scholes and Merton, 1973]\label{thm:blackscholes}
The strategy $a_t=\Phi(d_+(t,S_t))$, $c_t=e^{-rt}(C(t,S_t)-a_tS_t)$ is self-financing, its value is $V_t=C(t,S_t)$, and $V_t\to(S_T-K)^+$ as $t\to T$. An investor who starts with $C(0,S_0)$ can therefore
reproduce the payoff of the call exactly, whatever the path of the stock.
\end{theorem}

\begin{proof}
\emph{The formula as an expectation.} With $Z\sim N(0,1)$ and $t<T$,
\[
C(t,s)=e^{-r\tau}\,\mathbb E\bigl[(se^{\sigma\sqrt\tau Z+(r-\sigma^2/2)\tau}-K)^+\bigr].
\]
Indeed the payoff is positive exactly when $Z>-d_-$, and $Ke^{-r\tau}P(Z>-d_-)=Ke^{-r\tau}\Phi(d_-)$; completing the square,
\[
\mathbb E\bigl[e^{\sigma\sqrt\tau Z-\sigma^2\tau/2}\mathbf 1_{\{Z>-d_-\}}\bigr]=P(Z>-d_--\sigma\sqrt\tau)=\Phi(d_+).
\]
By dominated convergence this expression tends to $(s-K)^+$ as $t\to T$, locally uniformly in $s$, so $C$ is continuous on $[0,T]\times]0,\infty[$.

\emph{The partial differential equation.} Since $d_+^2-d_-^2=(d_+-d_-)(d_++d_-)=2(\log\frac sK+r\tau)$, we have $s\varphi(d_+)=Ke^{-r\tau}\varphi(d_-)$. As $d_+-d_-=\sigma\sqrt\tau$ does not depend on $s$,
\[
\partial_sC=\Phi(d_+)+s\varphi(d_+)\partial_sd_+-Ke^{-r\tau}\varphi(d_-)\partial_sd_-=\Phi(d_+),\qquad\partial_{ss}C=\frac{\varphi(d_+)}{s\sigma\sqrt\tau} ,
\]
and in the same way, using $\partial_t(d_+-d_-)=-\frac{\sigma}{2\sqrt\tau}$,
\[
\partial_tC=-\frac{s\varphi(d_+)\sigma}{2\sqrt\tau}-rKe^{-r\tau}\Phi(d_-).
\]
Therefore, on $[0,T[\times]0,\infty[$, where $C$ is $C^\infty$,
\begin{equation}\label{eq:bspde}
\partial_tC+rs\,\partial_sC+\frac{\sigma^2s^2}2\partial_{ss}C=rC.
\end{equation}

\emph{Replication.} $S$ stays in $]0,\infty[$, and It\^o's formula (in the $C^{1,2}$ form, on $[0,T-\varepsilon]$, after stopping when $S$ leaves $[1/k,k]$ and letting $k\to\infty$) gives, with
$\dd\langle S\rangle=\sigma^2S^2\dd t$ and \eqref{eq:bspde},
\[
\dd C(t,S_t)=\partial_sC\dd S_t+\Bigl(\partial_tC+\frac{\sigma^2S_t^2}2\partial_{ss}C\Bigr)\dd t=a_t\dd S_t+r\bigl(C(t,S_t)-a_tS_t\bigr)\dd t=a_t\dd S_t+c_t\,re^{rt}\dd t.
\]
So $V_t=a_tS_t+c_te^{rt}=C(t,S_t)$ satisfies the self-financing identity, and $V_t\to C(T,S_T)=(S_T-K)^+$ by continuity.
\end{proof}

\begin{remark}
If the call were sold at a price $p>C(0,S_0)$, one could sell it, run the strategy with $C(0,S_0)$, and keep $p-C(0,S_0)$ without risk; if $p<C(0,S_0)$, the reverse trade works. This \emph{no-arbitrage}
argument makes $C(0,S_0)$ the only rational price. The drift $\mu$ has disappeared: two investors who disagree about the expected return of the stock agree on the price of the option. The expectation
formula in the proof says that $C(0,S_0)$ is the discounted mean payoff for a stock with drift $r$ instead of $\mu$: by Girsanov's theorem (Theorem~\ref{thm:girsanov}) the price is an expectation under the
\emph{risk-neutral} probability, under which $e^{-rt}S_t$ is a martingale. This is the starting point of mathematical finance; see Karatzas and Shreve (1998) or Shreve (2004).
\end{remark}

\begin{think}
The strategy holds $\Phi(d_+)$ shares, a number between $0$ and $1$ which moves continuously. What happens to $a_t$ just before maturity, according to whether $S_t$ is above or below $K$? And why can a
strategy that only buys and sells the stock reproduce a payoff that is not linear in $S_T$?
\end{think}

\section{Bessel processes: which dimensions hit zero}

Let $B$ be a Brownian motion in $\R^d$ from $x$ and $Z=|B|^2$. It\^o's formula gives $Z_t=|x|^2+2\sum_i\int_0^tB^i\dd B^i+dt$. The martingale part is a time-changed Brownian motion run with the
speed $4Z$:

\begin{proposition}\label{prop:besqd}
Let $\beta_t=\sum_{i=1}^d\int_0^t\frac{B^i_s}{|B_s|}\mathbf 1_{\{B_s\ne0\}}\dd B^i_s+\int_0^t\mathbf 1_{\{B_s=0\}}\dd B^1_s$. Then $\beta$ is a Brownian motion and
\[
Z_t=|x|^2+dt+2\int_0^t\sqrt{Z_s}\dd\beta_s.
\]
\end{proposition}

\begin{proof}
$\beta$ is a continuous local martingale with $\langle\beta\rangle_t=\int_0^t(\mathbf 1_{\{B_s\ne0\}}+\mathbf 1_{\{B_s=0\}})\dd s=t$, a Brownian motion by L\'evy's characterisation. By Proposition~\ref{prop:itoprops}(b),
$\int\sqrt{Z}\dd\beta=\sum_i\int|B|\frac{B^i}{|B|}\mathbf 1_{\{B\ne0\}}\dd B^i=\sum_i\int B^i\dd B^i$, since $\sqrt Z=0$ on $\{B=0\}$.
\end{proof}

The equation $\dd Z=\delta\dd t+2\sqrt{Z}\dd\beta$ makes sense for every real $\delta\ge0$, and its coefficient $2\sqrt z$ is $\frac12$-H\"older, so by Theorem~\ref{thm:yamada} (applied with $\sigma(z)=2\sqrt{|z|}$)
it has at most one solution for a given $\beta$.

\begin{definition}\label{def:besq}
Let $\delta\ge0$ and $z\ge0$. A \emph{squared Bessel process of dimension $\delta$} from $z$ is a continuous adapted process $Z\ge0$ with $Z_t=z+\delta t+2\int_0^t\sqrt{Z_s}\dd\beta_s$ for a Brownian motion $\beta$.
$\sqrt Z$ is a \emph{Bessel process} of dimension $\delta$.
\end{definition}

For integer $\delta=d$, Proposition~\ref{prop:besqd} constructs one: the squared distance from the origin of $d$-dimensional Brownian motion. For non-integer $\delta$ a solution also exists (Revuz and Yor,
Chapter~XI); we shall not need it. The dimension decides whether the origin is reached, and the answer extends Corollary~\ref{cor:recurrence} to fractional dimensions.

\begin{theorem}\label{thm:besqzero}
Let $Z$ be a squared Bessel process of dimension $\delta>0$ from $z>0$, and $T_0=\inf\{t:Z_t=0\}$. If $\delta\ge2$, then $T_0=\infty$ almost surely; if $0<\delta<2$, then $T_0<\infty$ almost surely.
\end{theorem}

\begin{proof}
Let $0<\varepsilon<z<R$, $T_a=\inf\{t:Z_t=a\}$ and $\tau=T_\varepsilon\wedge T_R$. First, $\tau\le T_R<\infty$ almost surely: stopped at $T_R\wedge t$ the stochastic integral has an integrand bounded by $\sqrt R$, so
$\mathbb EZ_{T_R\wedge t}=z+\delta\,\mathbb E[T_R\wedge t]\le R$, and $\mathbb ET_R\le R/\delta$ by monotone convergence.

The generator of $Z$ acts on $C^2$ functions by $\mathcal Lf(z)=\delta f'(z)+2zf''(z)$, since $\dd\langle Z\rangle=4Z\dd t$. The \emph{scale function}
\[
s(z)=z^{1-\delta/2}\ (\delta\ne2),\qquad s(z)=\log z\ (\delta=2)
\]
satisfies $\mathcal Ls=0$ on $]0,\infty[$: for $p=1-\frac\delta2$, $\mathcal Lz^p=pz^{p-1}(\delta+2(p-1))=0$, and $\mathcal L\log z=\frac{\delta-2}z$. By It\^o's formula (for a $C^2$ function equal to $s$ on $[\varepsilon,R]$),
$s(Z_{t\wedge\tau})$ is a bounded martingale, so $s(z)=\mathbb Es(Z_\tau)$, that is
\[
P(T_\varepsilon<T_R)=\frac{s(R)-s(z)}{s(R)-s(\varepsilon)}.
\]
\emph{Case $\delta\ge2$.} As $\varepsilon\to0$, $s(\varepsilon)\to-\infty$ if $\delta=2$ and $s(\varepsilon)\to+\infty$ if $\delta>2$, so $P(T_\varepsilon<T_R)\to0$. Since $T_0\ge T_\varepsilon$, $P(T_0<T_R)=0$ for every $R$. A
continuous path is bounded on $[0,T_0]$ when $T_0<\infty$, so $\{T_0<\infty\}\subseteq\bigcup_R\{T_0<T_R\}$, a null set.

\emph{Case $\delta<2$.} Now $p>0$ and $s(\varepsilon)\to0$: $P(T_\varepsilon<T_R)\to1-(z/R)^p$. The events $\{T_\varepsilon<T_R\}$ decrease as $\varepsilon\downarrow0$; on their intersection $T_\varepsilon<T_R<\infty$ for all
$\varepsilon$, and $T_\varepsilon$ increases to a limit at which $Z=0$ by continuity, so $T_0\le T_R<\infty$. Hence $P(T_0<\infty)\ge1-(z/R)^p$ for every $R$, and $T_0<\infty$ almost surely.
\end{proof}

For $\delta=2$ this is again the fact that planar Brownian motion never hits a given point; for $\delta=1$, $\sqrt Z$ is $|x+B|$, which does. The threshold $2$ is a dimension, but nothing forces it to be an
integer. The next section meets a process that lives exactly at the threshold.

\section{Dyson Brownian motion}

Let $\mathcal H_n$ be the real vector space of Hermitian $n\times n$ matrices, of dimension $n^2$, with the inner product $\langle A,A'\rangle=\operatorname{tr}(AA')$. An orthonormal basis is formed by the
$e_je_j^*$, the $\frac1{\sqrt2}(e_je_l^*+e_le_j^*)$ and the $\frac{i}{\sqrt2}(e_je_l^*-e_le_j^*)$, $j<l$. A \emph{Hermitian Brownian motion} is $H_t=\sum_\alpha B^\alpha_tE_\alpha$, where $(E_\alpha)$ is an orthonormal
basis and the $B^\alpha$ are independent Brownian motions: the diagonal entries are real Brownian motions and the entries above the diagonal are $\frac1{\sqrt2}$ times complex Brownian motions. An orthogonal change
of basis maps an $n^2$-dimensional Brownian motion to another (Exercise~\ref{ex:rotation}), so the definition does not depend on the basis; in particular $UH_tU^*$ is again a Hermitian Brownian motion for a fixed
unitary $U$.

Write $\lambda_1(H)\le\dots\le\lambda_n(H)$ for the eigenvalues of $H$, and $\mathcal G$ for the set of matrices with $n$ distinct eigenvalues, which is open (it is $\{\mathrm{disc}\ne0\}$, where the
discriminant of the characteristic polynomial is a polynomial in the entries).

\begin{lemma}[eigenvalue perturbation]\label{lem:perturb}
The functions $\lambda_i$ are $C^\infty$ on $\mathcal G$. If $H\in\mathcal G$ has orthonormal eigenvectors $u_1,\dots,u_n$ and $E\in\mathcal H_n$, then
\[
D\lambda_i(H)[E]=u_i^*Eu_i,\qquad D^2\lambda_i(H)[E,E]=2\sum_{k\ne i}\frac{|u_k^*Eu_i|^2}{\lambda_i-\lambda_k}.
\]
\end{lemma}

\begin{proof}
Since $\lambda_i(UHU^*)=\lambda_i(H)$ for unitary $U$, we may assume $H=D=\mathrm{diag}(\lambda_1,\dots,\lambda_n)$ and $u_k=e_k$. Fix $i$ and consider, for $H'$ near $D$, $v\in e_i^\perp\subseteq\mathbb C^n$ and
$\mu\in\mathbb C$, the smooth map $G(v,\mu;H')=H'(e_i+v)-\mu(e_i+v)$. Its differential in $(v,\mu)$ at $(0,\lambda_i;D)$ is $(w,\nu)\mapsto(D-\lambda_i)w-\nu e_i$, a bijection from $e_i^\perp\times\mathbb C$ onto
$\mathbb C^n$ because $D-\lambda_i$ is invertible on $e_i^\perp$. The implicit function theorem gives smooth $v(H')$, $\mu_i(H')$ near $D$ with $G=0$: $\mu_i(H')$ is an eigenvalue of $H'$. Doing this for each $i$,
the $n$ eigenvalues $\mu_1(H'),\dots,\mu_n(H')$ are close to the distinct $\lambda_1<\dots<\lambda_n$, hence distinct, hence they are all the eigenvalues of $H'$, in increasing order: $\mu_i=\lambda_i$ near $D$.

For the derivatives, put $H'=D+\varepsilon E$, $\lambda(\varepsilon)=\lambda_i(D+\varepsilon E)$ and $v(\varepsilon)$ as above. The $i$-th coordinate of $G=0$ reads
$\lambda(\varepsilon)=\lambda_i+\varepsilon E_{ii}+\varepsilon\sum_{k\ne i}E_{ik}v_k(\varepsilon)$, and the $k$-th coordinate, $k\ne i$, reads $(\lambda(\varepsilon)-\lambda_k)v_k(\varepsilon)=\varepsilon(E(e_i+v(\varepsilon)))_k$, so that
$v_k(\varepsilon)=\varepsilon\frac{E_{ki}}{\lambda_i-\lambda_k}+O(\varepsilon^2)$. Hence
\[
\lambda(\varepsilon)=\lambda_i+\varepsilon E_{ii}+\varepsilon^2\sum_{k\ne i}\frac{E_{ik}E_{ki}}{\lambda_i-\lambda_k}+O(\varepsilon^3),
\]
and $E_{ik}E_{ki}=|E_{ki}|^2$ since $E$ is Hermitian. The first two derivatives of the smooth function $\lambda$ at $0$ are read off.
\end{proof}

\begin{figure}[ht]
\centering
\fig{ch06-dyson}
\caption{Left: the eigenvalues of $H_0+W_t$, with $H_0=\mathrm{diag}(-1,-\frac12,0,\frac12,1)$ and $H$ a Hermitian Brownian motion, on $[0,3]$. Right: five independent Brownian motions from the same points.
The eigenvalues repel and never cross; the independent paths cross freely.}
\label{fig:dyson}
\end{figure}

\begin{theorem}[Dyson, 1962; McKean, 1969]\label{thm:dyson}
Let $H_0\in\mathcal G$ and $H_t=H_0+W_t$ with $W$ a Hermitian Brownian motion. Almost surely, $H_t\in\mathcal G$ for all $t\ge0$: the eigenvalues never collide. The eigenvalues $\lambda_i(t)=\lambda_i(H_t)$ satisfy
\[
\lambda_i(t)=\lambda_i(0)+\beta_i(t)+\int_0^t\sum_{k\ne i}\frac{\mathrm ds}{\lambda_i(s)-\lambda_k(s)},\qquad i=1,\dots,n,
\]
where $\beta_1,\dots,\beta_n$ are independent Brownian motions.
\end{theorem}

\begin{proof}
\emph{Step 1: the generator.} Let $f\in C^2$ on the open set $\Omega=\{x\in\R^n:x_1<\dots<x_n\}$ and $F=f\circ\lambda$ on $\mathcal G$, which is $C^2$ by Lemma~\ref{lem:perturb}. Put $b_i(x)=\sum_{k\ne i}(x_i-x_k)^{-1}$. We claim that
for any orthonormal basis $(E_\alpha)$ of $\mathcal H_n$ and $H\in\mathcal G$, with $x=\lambda(H)$,
\begin{equation}\label{eq:dysongen}
\begin{gathered}
\sum_\alpha D\lambda_i(H)[E_\alpha]\,D\lambda_j(H)[E_\alpha]=\delta_{ij},\\
\frac12\sum_\alpha D^2F(H)[E_\alpha,E_\alpha]=\frac12\Delta f(x)+\sum_ib_i(x)\partial_if(x).
\end{gathered}
\end{equation}
Both sums are unchanged by an orthogonal change of basis (they are the trace of a quadratic form, and a sum of products of coordinates of two vectors), so we may use the basis $UE'_\alpha U^*$, where
$H=U\mathrm{diag}(x)U^*$ and $(E'_\alpha)$ is the basis listed above. By Lemma~\ref{lem:perturb}, $D\lambda_i[Ue_je_j^*U^*]=\delta_{ij}$ and $D^2\lambda_i[Ue_je_j^*U^*,\cdot]=0$; for the two basis elements attached to
$j<l$, $D\lambda_i=0$ (their diagonal entries vanish) and $D^2\lambda_i=2\cdot\frac12\bigl(\frac{\delta_{ij}}{x_j-x_l}+\frac{\delta_{il}}{x_l-x_j}\bigr)$ for each. The first identity follows, and summing,
$\sum_\alpha D^2\lambda_i[E_\alpha,E_\alpha]=2b_i(x)$. The chain rule $D^2F[E,E]=\sum_{i,j}\partial_{ij}f\,D\lambda_i[E]D\lambda_j[E]+\sum_i\partial_if\,D^2\lambda_i[E,E]$ gives the second identity.

\emph{Step 2: It\^o's formula.} For $\varepsilon,R>0$ let $\mathcal K_{\varepsilon,R}=\{H:\|H\|\le R,\ \lambda_{i+1}(H)-\lambda_i(H)\ge\varepsilon\ \forall i\}$, a compact subset of $\mathcal G$, and $\tau_{\varepsilon,R}$ the exit time of $H_t$ from it.
Multiplying $F$ by a smooth function equal to $1$ on $\mathcal K_{\varepsilon,R}$ and supported in $\mathcal G$, It\^o's formula in the $n^2$ coordinates $B^\alpha$ gives, for $t\le\tau_{\varepsilon,R}$, by \eqref{eq:dysongen},
\[
f(\lambda(t))=f(\lambda(0))+\sum_i\int_0^t\partial_if(\lambda(s))\dd\beta_i(s)+\int_0^t\Bigl(\frac12\Delta f+\sum_ib_i\partial_if\Bigr)(\lambda(s))\dd s,
\]
where $\beta_i(t)=\sum_\alpha\int_0^tD\lambda_i(H_s)[E_\alpha]\mathbf 1_{\{H_s\in\mathcal G\}}\dd B^\alpha_s$ (an integrand bounded by $\|E_\alpha\|\le1$).

\emph{Step 3: McKean's Lyapunov function.} Take $f=\psi$, with
\[
\psi(x)=\sum_ix_i^2-\sum_{i<j}\log(x_j-x_i).
\]
For $g(x)=\sum_{i<j}\log(x_j-x_i)$ we have $\partial_ig=b_i$, so $\frac12\Delta g=-\frac12S$ with $S=\sum_{i\ne k}(x_i-x_k)^{-2}$, and
\[
\sum_ib_i\partial_ig=\sum_ib_i^2=S+\sum_{i,k,l\ \text{distinct}}\frac1{(x_i-x_k)(x_i-x_l)}=S,
\]
because the sum over distinct triples vanishes: grouping the six orderings of $\{i,k,l\}$, it is a sum of terms $\frac{(x_k-x_l)+(x_l-x_i)+(x_i-x_k)}{(x_i-x_k)(x_k-x_l)(x_l-x_i)}=0$. So the drift of $g$ is $\frac12S$.
For $\sum x_i^2$ the drift is $n+2\sum_ix_ib_i=n+\sum_{i\ne k}\frac{2x_i}{x_i-x_k}=n+\sum_{i<k}2=n^2$. Hence $\frac12\Delta\psi+\sum b_i\partial_i\psi=n^2-\frac12S\le n^2$.

On the other hand $\psi$ is bounded below. Since $\log y\le cy^2+C_c$ for every $c>0$ and $(x_j-x_i)^2\le2x_i^2+2x_j^2$, choosing $c$ small gives $\sum_{i<j}\log(x_j-x_i)\le\frac12\sum_ix_i^2+C$, with $C$ depending on $n$ only.
Hence $\psi\ge-C$, and more precisely, if some gap $x_{i+1}-x_i$ equals $\varepsilon$, then, isolating that term, $\psi(x)\ge\log\frac1\varepsilon-C'$.

The stochastic integral stopped at $\tau_{\varepsilon,R}$ has bounded integrands, so taking expectations,
$\mathbb E\psi(\lambda(t\wedge\tau_{\varepsilon,R}))\le\psi(\lambda(0))+n^2t$. Let $\tau_\varepsilon=\inf\{t:\min_i(\lambda_{i+1}-\lambda_i)(t)\le\varepsilon\}$; since $\|H_t\|$ is bounded on bounded intervals, $\tau_{\varepsilon,R}=\tau_\varepsilon$ for $R$
large (depending on $\omega$ and $t$), and Fatou's lemma applied to $\psi+C\ge0$ gives $\mathbb E[\psi(\lambda(t\wedge\tau_\varepsilon))+C]\le\psi(\lambda(0))+n^2t+C$. On $\{\tau_\varepsilon\le t\}$ a gap equals $\varepsilon$ at time $\tau_\varepsilon$, so
\[
P(\tau_\varepsilon\le t)\Bigl(\log\frac1\varepsilon-C'+C\Bigr)\le\psi(\lambda(0))+n^2t+C.
\]
A collision before $t$ implies $\tau_\varepsilon\le t$ for every $\varepsilon$, so its probability is $0$. This holds for every $t$: almost surely $H_t\in\mathcal G$ for all $t$.

\emph{Step 4: the equation.} The $\beta_i$ are continuous local martingales with $\langle\beta_i,\beta_j\rangle_t=\int_0^t\sum_\alpha D\lambda_i[E_\alpha]D\lambda_j[E_\alpha]\dd s=\delta_{ij}t$ by \eqref{eq:dysongen}, so they are independent
Brownian motions (Theorem~\ref{thm:levychar}). Step 2 with $f(x)=x_i$, valid up to $\tau_{\varepsilon,R}$ for all $\varepsilon,R$, hence for all $t$ by Step 3, is the stated equation.
\end{proof}

\begin{beauty}
Nothing in the definition of $H_t$ prevents two eigenvalues from meeting: the matrices with a repeated eigenvalue are a perfectly good subset of $\mathcal H_n$. But it has codimension $3$, and the
eigenvalues feel it as a force $\frac1{\lambda_i-\lambda_k}$, exactly the force that makes a three-dimensional Brownian motion avoid the origin. For $n=2$ this is literal: the gap
$\lambda_2-\lambda_1$, divided by $\sqrt2$, is a Bessel process of dimension $3$ (Exercise~\ref{ex:twobytwo}). The repulsion that random matrix theory finds between the zeros of the zeta function
(Chapter~11) is already here.
\end{beauty}

\begin{remark}
Dyson also considered real symmetric and quaternionic matrices; the eigenvalues then satisfy the same equation with the Brownian terms multiplied by $\sqrt{2/\beta}$, where $\beta=1,2,4$ is Dyson's index (in a suitable normalisation), and
they never collide for any $\beta\ge1$ (Anderson, Guionnet and Zeitouni, Section~4.3). The theorem above starts from $H_0$ with distinct eigenvalues. The entrance from a multiple eigenvalue, in particular from $H_0=0$, is not proved here, and neither is the claim that $\lambda(1)$ has the joint density proportional to
$\prod_{i<j}(x_j-x_i)^2e^{-\sum x_i^2/2}$. Chapter~10 studies the GOE semicircle, the Marchenko--Pastur law and a Wigner spike; it does not derive that Gaussian unitary density.
\end{remark}

\section{Solving equations numerically: Euler--Maruyama}

The simplest scheme freezes the coefficients on each step. For a step $h>0$ let $\eta(s)=h\lfloor s/h\rfloor$ and define $X^h$ by $X^h_0=x$ and
\[
X^h_t=X^h_{\eta(t)}+b(X^h_{\eta(t)})(t-\eta(t))+\sigma(X^h_{\eta(t)})(B_t-B_{\eta(t)}) ,
\]
so that $X^h$ at the grid times is computed with one Gaussian vector per step, and $X^h_t=x+\int_0^tb(X^h_{\eta(s)})\dd s+\int_0^t\sigma(X^h_{\eta(s)})\dd B_s$. This is Maruyama's 1955 scheme, Euler's method with noise.

\begin{theorem}[strong order $\frac12$]\label{thm:eulermaruyama}
Under the hypotheses of Theorem~\ref{thm:sdeexist}, let $X$ be the solution from $x$. For every $T$ there is $C$, depending on $T,K,b(0),\sigma(0)$ and $|x|$, with
$\mathbb E\sup_{t\le T}|X_t-X^h_t|^2\le Ch$ for all $h\le1$.
\end{theorem}

\begin{proof}
As for Theorem~\ref{thm:sdeexist}, Gronwall applied to $\mathbb E\sup_{s\le t}|X^h_s|^2$ gives $\mathbb E\sup_{t\le T}|X^h_t|^2\le C_1$, uniformly in $h$. Then for every $s\le T$,
\[
\mathbb E|X^h_s-X^h_{\eta(s)}|^2\le2h^2\,\mathbb E|b(X^h_{\eta(s)})|^2+2h\,\mathbb E|\sigma(X^h_{\eta(s)})|^2\le C_2h.
\]
Let $e_t=X_t-X^h_t=\int_0^t(b(X_s)-b(X^h_{\eta(s)}))\dd s+\int_0^t(\sigma(X_s)-\sigma(X^h_{\eta(s)}))\dd B_s$. By Cauchy--Schwarz, Doob's inequality, \eqref{eq:isomvec} and
$|X_s-X^h_{\eta(s)}|\le|e_s|+|X^h_s-X^h_{\eta(s)}|$, the function $\varphi(t)=\mathbb E\sup_{s\le t}|e_s|^2$ satisfies
with $L=2(T+4)K^2$,
\[
\varphi(t)\le L\int_0^t\mathbb E|X_s-X^h_{\eta(s)}|^2\dd s\le2L\int_0^t\bigl(\varphi(s)+C_2h\bigr)\dd s ,
\]
and Gronwall gives $\varphi(T)\le2LTC_2h\,e^{2LT}$.
\end{proof}

The rate $\sqrt h$ for the error is the size of one Brownian increment, and it is sharp in general (Exercise~\ref{ex:emgbm}). Adding the next term of the It\^o--Taylor expansion,
$\frac12\sigma\sigma'(X^h_{\eta(t)})((B_t-B_{\eta(t)})^2-(t-\eta(t)))$ in dimension one, gives Milstein's scheme (1974), of strong order $1$; and for expectations $\mathbb Ef(X_T)$ of smooth $f$ the Euler scheme
already has \emph{weak} order $1$ (Kloeden and Platen, 1992). Both convergence orders can be measured numerically.

\section{Exercises}

\begin{exercise}[$\star$]\label{ex:nonunique}
Show that $X=0$ and $X=B^3$ both solve $\dd X=3X^{1/3}\dd t+3X^{2/3}\dd B$, $X_0=0$ (real cube roots). Which hypothesis of Theorem~\ref{thm:yamada} fails?
\end{exercise}

\begin{exercise}[$\star$]\label{ex:gbm}
Let $S$ be the geometric Brownian motion of Proposition~\ref{prop:gbm}. Show that $\mathbb ES_t=se^{\mu t}$, and that $S_t\to0$ almost surely if $\mu<\frac{\sigma^2}2$. What happens to a gambler who, at each instant,
invests a fixed fraction of a fortune growing at mean rate $\mu$ with $0<\mu<\frac{\sigma^2}2$?
\end{exercise}

\begin{exercise}[$\star$]\label{ex:putcall}
In the Black--Scholes model, a \emph{put} pays $(K-S_T)^+$. Show that $(S_T-K)^+-(K-S_T)^+=S_T-K$ and deduce that the price of the put is $C(t,s)-s+Ke^{-r(T-t)}$. Show that $C$ increases with the
volatility $\sigma$: $\partial_\sigma C=s\varphi(d_+)\sqrt{T-t}>0$.
\end{exercise}

\begin{exercise}[$\star\star$]\label{ex:bridgesde}
Let $X_t=(1-t)\int_0^t\frac{\mathrm dB_s}{1-s}$ for $t<1$. Show that $X$ solves $\dd X=-\frac{X}{1-t}\dd t+\dd B$, that it is a centred Gaussian process with covariance $s(1-t)$ for $s\le t<1$ (a Brownian bridge),
and that $X_t\to0$ almost surely as $t\to1$.
\end{exercise}

\begin{exercise}[$\star\star$]\label{ex:langevin}
Let $X$ be the stationary Ornstein--Uhlenbeck process ($\theta,\sigma>0$, $X_0\sim N(0,\frac{\sigma^2}{2\theta})$ independent of $B$) and $Y_t=\int_0^tX_s\dd s$. Show that $\theta Y_t=X_0-X_t+\sigma B_t$, and deduce that
$Y_{ct}/\sqrt c$ converges in law to $\frac\sigma\theta B_t$ as $c\to\infty$, and that $\mathbb EY_t^2\sim\frac{\sigma^2}{\theta^2}t$.
\end{exercise}

\begin{exercise}[$\star\star$]\label{ex:bessel3}
Let $Z$ be the squared Bessel process of dimension $\delta>1$ from $z>0$, and suppose $T_0=\infty$ (which holds for $\delta\ge2$). Show that $R=\sqrt Z$ satisfies $R_t=\sqrt z+\beta_t+\frac{\delta-1}2\int_0^t\frac{\mathrm ds}{R_s}$.
For $\delta=3$, show that $\mathbb E\int_0^t\frac{\mathrm ds}{R_s}<\infty$ and that $R_t\to\infty$ almost surely.
\end{exercise}

\begin{exercise}[$\star\star$]\label{ex:twobytwo}
For $n=2$, write $H_t=H_0+W_t$ as $\frac{\operatorname{tr}H_t}2I+\begin{pmatrix}a_t&c_t\\\bar c_t&-a_t\end{pmatrix}$. Show that $\lambda_2-\lambda_1=2\sqrt{a^2+|c|^2}$, that $(a,\operatorname{Re}c,\operatorname{Im}c)$ is a
three-dimensional Brownian motion with variance $\frac12$ per coordinate, and deduce that $(\lambda_2-\lambda_1)/\sqrt2$ is a Bessel process of dimension $3$. Check this against Theorem~\ref{thm:dyson}.
\end{exercise}

\begin{exercise}[$\star\star$]\label{ex:emgbm}
Let $d=m=1$, $b(x)=0$, $\sigma(x)=x$ and $x=1$, so that $X_t=e^{B_t-t/2}$, and let $X^h$ be the Euler--Maruyama approximation. Show that $X^h_1=\prod_{k=1}^{1/h}(1+\Delta_kB)$ ($1/h$ an integer), and that
$\mathbb E|X_1-X^h_1|^2\ge ch$ for some $c>0$ and all small $h$, so that the order $\frac12$ cannot be improved for this scheme.
\end{exercise}

\begin{exercise}[$\star\star\star$]\label{ex:sdegenerator}
Under the hypotheses of Theorem~\ref{thm:sdeexist}, let $f\in C^2$ with compact support and $\mathcal Lf=b\cdot\nabla f+\frac12\sum_{i,j}(\sigma\sigma^{\mathsf T})_{ij}\partial_{ij}f$. Show that
$f(X^x_t)-f(x)-\int_0^t\mathcal Lf(X^x_s)\dd s$ is a martingale, and that $\frac1t(P_tf(x)-f(x))\to\mathcal Lf(x)$ as $t\to0$, uniformly in $x$ if $b$ and $\sigma$ are bounded.
\end{exercise}

\section{Hints}

\begin{hint}{ex:nonunique}
It\^o's formula for $B^3$. Is $x\mapsto3x^{1/3}$ Lipschitz?
\end{hint}

\begin{hint}{ex:gbm}
$\mathbb Ee^{\sigma B_t}=e^{\sigma^2t/2}$, and $B_t/t\to0$.
\end{hint}

\begin{hint}{ex:putcall}
Both sides of the parity are replicated by static portfolios. For $\partial_\sigma C$, use $s\varphi(d_+)=Ke^{-r\tau}\varphi(d_-)$ as in the proof of Theorem~\ref{thm:blackscholes}.
\end{hint}

\begin{hint}{ex:bridgesde}
It\^o's formula for the product of $(1-t)$ and a martingale. For the limit, write $M_t=\int_0^t\frac{\mathrm dB_s}{1-s}$ as a time-changed Brownian motion with clock $\frac t{1-t}$.
\end{hint}

\begin{hint}{ex:langevin}
Integrate the equation. $X_0$ and $X_t$ stay bounded in $L^2$.
\end{hint}

\begin{hint}{ex:bessel3}
It\^o's formula for $\sqrt z$, which is smooth on $]0,\infty[$. For $\delta=3$, compute $\mathbb E[1/R_s]$ from $R_s=|x+B_s|$ with $B$ three-dimensional.
\end{hint}

\begin{hint}{ex:twobytwo}
The traceless part is $\sum$ of three orthogonal basis elements with Brownian coefficients. Then use Proposition~\ref{prop:besqd} and Exercise~\ref{ex:bessel3}.
\end{hint}

\begin{hint}{ex:emgbm}
Compare $\log X^h_1=\sum\log(1+\Delta_kB)$ with $B_1-\frac12$, using $\log(1+y)=y-\frac{y^2}2+\frac{y^3}3+O(y^4)$; the cubic terms have variance of order $h^2$ each.
\end{hint}

\begin{hint}{ex:sdegenerator}
It\^o's formula, then take expectations; for the limit, $\mathcal Lf$ is bounded and continuous, and $X^x_s\to x$ as $s\to0$.
\end{hint}

\section{Solutions}

\begin{solution}{ex:nonunique}
\idea Cubing a Brownian motion produces an It\^o drift $3B\dd t$; expressed in terms of $X=B^3$, both coefficients become powers of $X$, and the zero solution sits alongside.
\writeup $\dd(B^3)=3B^2\dd B+3B\dd t$, and $B^2=X^{2/3}$, $B=X^{1/3}$ with real cube roots. The zero function is also a solution. The diffusion coefficient $3x^{2/3}$ is $\frac12$-H\"older near $0$, so the
hypothesis on $\sigma$ in Theorem~\ref{thm:yamada} holds locally; it is the drift $3x^{1/3}$ that is not Lipschitz at $0$, exactly as in the deterministic example $x'=3x^{2/3}$, which has the solutions $0$ and $t^3$.
\end{solution}

\begin{solution}{ex:gbm}
\idea The mean is carried by rare paths on which $B$ is large; the typical path is governed by the exponent $\sigma B_t+(\mu-\frac{\sigma^2}2)t$, whose drift wins over $B_t$.
\writeup $\mathbb ES_t=se^{(\mu-\sigma^2/2)t}\mathbb Ee^{\sigma B_t}=se^{\mu t}$. Since $B_t/t=X_{1/t}\to0$ almost surely, where $X_u=uB_{1/u}$ is the time-inverted Brownian motion of Chapter~2,
$\frac1t\log S_t\to\mu-\frac{\sigma^2}2<0$, so $S_t\to0$. The gambler's expected fortune grows exponentially, and yet he is ruined almost surely: the long-run growth rate of a typical fortune is
$\mu-\frac{\sigma^2}2$, the quantity that Kelly's criterion (1956) maximises.
\end{solution}

\begin{solution}{ex:putcall}
\idea A call minus a put is a forward contract, worth today the stock minus the discounted strike; and more volatility makes the option's convex payoff more valuable.
\writeup The identity $y^+-(-y)^+=y$ with $y=S_T-K$. Holding one share and borrowing $Ke^{-rT}$ is a self-financing strategy (constant holdings) worth $S_t-Ke^{-r(T-t)}$ at time $t$ and $S_T-K$ at $T$; the strategy of
Theorem~\ref{thm:blackscholes} minus this one replicates $(S_T-K)^+-(S_T-K)=(K-S_T)^+$, starting from $C(t,s)-s+Ke^{-r(T-t)}$, which is therefore the price of the put. Differentiating, $\partial_\sigma C=s\varphi(d_+)\partial_\sigma d_+-Ke^{-r\tau}\varphi(d_-)\partial_\sigma d_-
=s\varphi(d_+)\partial_\sigma(d_+-d_-)=s\varphi(d_+)\sqrt\tau$.
\end{solution}

\begin{solution}{ex:bridgesde}
\idea The factor $(1-t)$ crushes the process to $0$ at time $1$, while the integral $\int\frac{\mathrm dB}{1-s}$ explodes only like $\sqrt{t/(1-t)}$.
\writeup Let $M_t=\int_0^t\frac{\mathrm dB_s}{1-s}$. It\^o's formula for $f(t,m)=(1-t)m$: $\dd X=-M\dd t+(1-t)\dd M=-\frac X{1-t}\dd t+\dd B$. $X$ is centred Gaussian (Wiener integrals), and for $s\le t$,
$\mathbb EX_sX_t=(1-s)(1-t)\int_0^s\frac{\mathrm du}{(1-u)^2}=(1-s)(1-t)\frac s{1-s}=s(1-t)$, the covariance of the Brownian bridge. $M$ is a continuous martingale with $\langle M\rangle_t=\frac t{1-t}$, so
$M_t=\gamma_{t/(1-t)}$ for a Brownian motion $\gamma$ (Exercise~\ref{ex:wienerito}(c) or Theorem~\ref{thm:dds}), and $X_t=(1-t)\gamma_{u}$ with $u=\frac t{1-t}\to\infty$, $1-t=\frac1{1+u}$; since $\gamma_u/u\to0$
almost surely (time inversion again), $X_t\to0$.
\end{solution}

\begin{solution}{ex:langevin}
\idea Integrating Langevin's equation, the friction term is $\theta$ times the displacement, and the force term is a Brownian motion; the velocity itself stays bounded and disappears at large scales.
\writeup $X_t=X_0-\theta\int_0^tX_s\dd s+\sigma B_t$, which is the identity. Then $Y_{ct}/\sqrt c=\frac\sigma\theta B_{ct}/\sqrt c+\frac{X_0-X_{ct}}{\theta\sqrt c}$; the first term has the law of $\frac\sigma\theta B_t$ and the second
tends to $0$ in $L^2$, since $\mathbb EX_u^2=\frac{\sigma^2}{2\theta}$ for all $u$. For the variance, $\mathbb EY_t^2=\theta^{-2}\mathbb E(X_0-X_t+\sigma B_t)^2=\frac{\sigma^2}{\theta^2}t+O(\sqrt t)$ by Cauchy--Schwarz. Einstein's
diffusion coefficient appears as $\frac{\sigma^2}{2\theta^2}$.
\end{solution}

\begin{solution}{ex:bessel3}
\idea On $]0,\infty[$ the square root is smooth and It\^o's formula turns the constant drift $\delta$ of $Z$ into a drift $\frac{\delta-1}{2R}$ that pushes away from $0$.
\writeup For $g(z)=\sqrt z$, $g'=\frac1{2\sqrt z}$, $g''=-\frac1{4z^{3/2}}$, and It\^o's formula, valid on $]0,\infty[$ (stop when $Z$ leaves $[1/k,k]$), gives
$\dd R=\frac1{2R}(\delta\dd t+2R\dd\beta)-\frac18\frac{4R^2}{R^3}\dd t=\dd\beta+\frac{\delta-1}{2R}\dd t$. For $\delta=3$ we may take $R_s=|x+B_s|$ with $|x|^2=z$ (Proposition~\ref{prop:besqd} and pathwise uniqueness).
Given $|B_s|=\rho$, the point $x+B_s$ is uniform on the sphere of centre $x$ and radius $\rho$ (Lemma~\ref{lem:uniformsphere}), and the mean of $\frac1{|y|}$ on that sphere is $\frac1{\max(|x|,\rho)}$: for $\rho<|x|$ by the
mean value property of the harmonic function $\frac1{|y|}$ (Lemma~\ref{lem:meanvalue}); for $\rho>|x|$ because the mean of $\frac1{|x-z|}$ over $|z|=\rho$ is a radial harmonic function of $x$ in the ball $|x|<\rho$, hence
equal to its value $\frac1\rho$ at the centre. So $\mathbb E\frac1{|x+B_s|}\le\mathbb E\frac1{|B_s|}=\sqrt{\frac2{\pi s}}$, which is integrable on $[0,t]$. Finally $R_t\to\infty$ is the transience of three-dimensional Brownian motion (Corollary~\ref{cor:recurrence}).
\end{solution}

\begin{solution}{ex:twobytwo}
\idea A $2\times2$ Hermitian matrix minus its trace is a vector in a three-dimensional space, its eigenvalues are $\pm$ its length, and the length of a three-dimensional Brownian motion is a Bessel(3) process.
\writeup The traceless matrix $A=\begin{pmatrix}a&c\\\bar c&-a\end{pmatrix}$ has $A^2=(a^2+|c|^2)I$, so its eigenvalues are $\pm\sqrt{a^2+|c|^2}$ and the gap is $2\sqrt{a^2+|c|^2}$. In the orthonormal basis of the text,
$a_t=a_0+\frac12(B^{11}_t-B^{22}_t)$, $\operatorname{Re}c_t=\operatorname{Re}c_0+\frac1{\sqrt2}B^{12}_t$, $\operatorname{Im}c_t=\operatorname{Im}c_0+\frac1{\sqrt2}B'^{12}_t$,
three independent Brownian motions each of variance $\frac12$ per unit time. So $\sqrt2(a,\operatorname{Re}c,\operatorname{Im}c)$ is a standard three-dimensional Brownian motion $Y$, the gap is $2|Y|/\sqrt2=\sqrt2|Y|$, and
$(\lambda_2-\lambda_1)/\sqrt2=|Y|$, a Bessel process of dimension $3$. Check: Theorem~\ref{thm:dyson} gives $\dd(\lambda_2-\lambda_1)=\dd(\beta_2-\beta_1)+\frac2{\lambda_2-\lambda_1}\dd t$, so $\rho=(\lambda_2-\lambda_1)/\sqrt2$ satisfies
$\dd\rho=\dd\gamma+\frac1\rho\dd t$ with $\gamma=(\beta_2-\beta_1)/\sqrt2$ a Brownian motion: the equation of Exercise~\ref{ex:bessel3} with $\delta=3$.
\end{solution}

\begin{solution}{ex:emgbm}
\idea The scheme multiplies by $1+\Delta B$ where the solution multiplies by $e^{\Delta B-h/2}$; the logarithms differ by about $\frac{(\Delta B)^2-h}2$ and $\frac{(\Delta B)^3}3$, and these errors, independent from step
to step, add up to a fluctuation of order $\sqrt h$.
\writeup With $b=0$ and $\sigma(x)=x$, $X^h_{(k+1)h}=X^h_{kh}(1+\Delta_{k+1}B)$. On the event $A_h=\{\max_k|\Delta_kB|\le\frac12\}$, whose probability tends to $1$, $\log X^h_1-\log X_1=\sum_k\bigl(\log(1+\Delta_kB)-\Delta_kB+\frac h2\bigr)
=-\sum_k\frac{(\Delta_kB)^2-h}2+\sum_k\frac{(\Delta_kB)^3}3+R$, with $|R|\le C\sum(\Delta_kB)^4$, $\mathbb E|R|\le C'h$. The first sum has variance $\frac1h\cdot\frac{2h^2}4=\frac h2$ and the second $\frac1h\cdot\frac{15h^3}9$; so
$D_h=(\log X^h_1-\log X_1)/\sqrt h$ converges in law to $N(0,\frac12)$ (central limit theorem for the first sum). Then $|X_1-X^h_1|=X_1|1-e^{\sqrt hD_h}|$, and on $A_h\cap\{X_1\ge\frac12\}\cap\{|D_h|\in[\frac12,2]\}$, an event of probability
at least some $p>0$ for small $h$, $|X_1-X^h_1|\ge\frac12\cdot\frac{\sqrt h}2\cdot\frac12$. This event has probability bounded below for small $h$: $P(X_1\ge\frac12)=P(B_1\ge\frac12-\log2)>0.57$, $P(|D_h|\in[\frac12,2])\to P(|N(0,\frac12)|\in[\frac12,2])>0.47$ and $P(A_h)\to1$,
so their intersection has probability at least $p=0.03$ for small $h$, by $P(E\cap F)\ge P(E)+P(F)-1$. Hence $\mathbb E|X_1-X^h_1|^2\ge p\,h/64$.
\end{solution}

\begin{solution}{ex:sdegenerator}
\idea It\^o's formula splits $f(X_t)$ into a martingale and the integral of $\mathcal Lf$; at small times the integral is $t\,\mathcal Lf(x)$ to first order.
\writeup It\^o's formula, with $\dd\langle X^i,X^j\rangle=(\sigma\sigma^{\mathsf T})_{ij}(X)\dd t$, gives $f(X_t)=f(x)+\int_0^t\nabla f(X)\cdot\sigma(X)\dd B+\int_0^t\mathcal Lf(X)\dd s$. The integrand $\nabla f\cdot\sigma$ is bounded
($f$ has compact support and $\sigma$ is continuous), so the stochastic integral is a martingale. Taking expectations, $\frac1t(P_tf(x)-f(x))-\mathcal Lf(x)=\frac1t\int_0^t\mathbb E[\mathcal Lf(X^x_s)-\mathcal Lf(x)]\dd s$. $\mathcal Lf$ is
bounded and continuous, and $X^x_s\to x$ almost surely as $s\to0$, so the right side tends to $0$ by dominated convergence. If $b,\sigma$ are bounded, $\mathcal Lf$ is uniformly continuous and
$\mathbb E|X^x_s-x|^2\le2\|b\|^2s^2+2\|\sigma\|^2s$ uniformly in $x$, which gives uniformity. Chapter~7 starts from this.
\end{solution}

\section*{Sources and further reading}

\begin{itemize}
\item \emph{History.} P.~Langevin, ``Sur la th\'eorie du mouvement brownien'', \emph{Comptes rendus de l'Acad\'emie des sciences} 146 (1908), 530--533; English translation by A.~Gythiel, with an introduction by D.~S.~Lemons,
\emph{American Journal of Physics} 65 (1997), 1079--1081. G.~E.~Uhlenbeck and L.~S.~Ornstein, ``On the theory of the Brownian motion'', \emph{Physical Review} 36 (1930), 823--841. K.~It\^o, \emph{On Stochastic Differential
Equations}, Memoirs of the American Mathematical Society 4, 1951. G.~Maruyama, ``Continuous Markov processes and stochastic equations'', \emph{Rendiconti del Circolo Matematico di Palermo} 4 (1955), 48--90.
F.~Black and M.~Scholes, ``The pricing of options and corporate liabilities'', \emph{Journal of Political Economy} 81 (1973), 637--654. R.~C.~Merton, ``Theory of rational option pricing'', \emph{Bell Journal of
Economics and Management Science} 4 (1973), 141--183. F.~Black, ``How we came up with the option formula'', \emph{Journal of Portfolio Management} 15 (1989), 4--8. F.~J.~Dyson, ``A Brownian-motion model for the
eigenvalues of a random matrix'', \emph{Journal of Mathematical Physics} 3 (1962), 1191--1198. H.~P.~McKean, \emph{Stochastic Integrals}, Academic Press, 1969. T.~Yamada and S.~Watanabe, ``On the uniqueness of
solutions of stochastic differential equations'', \emph{Journal of Mathematics of Kyoto University} 11 (1971), 155--167. J.~L.~Kelly, ``A new interpretation of information rate'', \emph{Bell System Technical Journal} 35 (1956), 917--926. G.~N.~Milstein, ``Approximate integration of stochastic differential equations'', \emph{Theory of
Probability and its Applications} 19 (1974), 557--562.
\item D.~Revuz and M.~Yor, \emph{Continuous Martingales and Brownian Motion}, 3rd edition, Springer, 1999, Chapters~IX and~XI. Strong and weak solutions, Yamada--Watanabe, Bessel processes in every dimension.
\item G.~W.~Anderson, A.~Guionnet and O.~Zeitouni, \emph{An Introduction to Random Matrices}, Cambridge University Press, 2010, Section~4.3. Dyson Brownian motion for $\beta=1,2,4$ and McKean's argument.
\item I.~Karatzas and S.~E.~Shreve, \emph{Methods of Mathematical Finance}, Springer, 1998; S.~E.~Shreve, \emph{Stochastic Calculus for Finance II}, Springer, 2004. Black--Scholes and risk-neutral pricing.
\item P.~E.~Kloeden and E.~Platen, \emph{Numerical Solution of Stochastic Differential Equations}, Springer, 1992. Euler--Maruyama, Milstein and higher-order schemes, strong and weak orders.
\end{itemize}
